\section{The canonical good model and the main theorem}
\label{sec:good-model-interface}

The geometric argument has used the existence of a good model of
$X$ through its sections and its effective exceptional comparison.
We now supply exactly that model. The following theorem is
\cite[Corollary~11.2]{LogAbundance2026}; its proof belongs to that
companion article.

\begin{theorem}[Smooth canonical good models]
\label{thm:canonical-good-models}
Let $T$ be a smooth connected projective complex variety with
pseudo-effective $K_T$. Then there exist a normal projective
$\Q$-factorial klt variety $V$ and a $K_T$-negative birational
contraction $T\dashrightarrow V$ such that $K_V$ is $\Q$-Cartier
and semiample. On a smooth projective common resolution
$T\xleftarrow{p}W\xrightarrow{q}V$, compatible canonical divisors
satisfy
\[
 p^*K_T=q^*K_V+E,\qquad E\geq0,\qquad q_*E=0.
\]
\end{theorem}

Thus the input provides both the actual globally generated multiple
used for sections and the exceptional equality used for intersections.

\begin{proof}[Proof of Theorem~\ref{thm:schnell}]
If $Y$ is a point, then its geometric generic fiber is $X$, and
the conclusion is the hypothesis $\kappa(F)=0$.
For a positive-dimensional base, a positive multiple of the ample
divisor $H$ has an effective representative. Its pullback shows that $f^*H$ is
pseudo-effective. Adding it to the original numerical hypothesis gives
\[
 K_X=\frac1{m_0}\bigl((m_0K_X-f^*H)+f^*H\bigr)\in\PE(X).
\]
The variety $X$ is smooth, connected, projective and complex, so
Theorem~\ref{thm:canonical-good-models} applies. Its output is the
projective good klt minimal model required by
Theorem~\ref{thm:good-model-case}. Applying that theorem to the
original $f$, $H$ and $m_0$ gives $\kappa(X)=\dim Y$.
\end{proof}

Throughout the proof the assumed comparison is numerical. The
pluricanonical sections are obtained from the semiample model through
\eqref{eq:section-transfer}; they are not assumed as a reformulation
of pseudo-effectivity.
